\documentclass[UTF8,a4paper,10pt,fontset=none]{ctexart}
\usepackage[margin=18mm,headheight=15pt]{geometry}
\usepackage{amsmath,amssymb,booktabs,tabularx,array,xcolor,fancyhdr,hyperref,enumitem}
\setmainfont{Arial}\setCJKmainfont{Microsoft YaHei}\setCJKmonofont{Microsoft YaHei}
\definecolor{ink}{HTML}{18344C}
\hypersetup{unicode=true,colorlinks=true,urlcolor=ink,linkcolor=ink,pdftitle={YunTAPR-Net 极端降水诊断一次执行准备与限额},pdfauthor={YunTAPR-Net Research}}
\pagestyle{fancy}\fancyhf{}\lhead{\small YunTAPR-Net}\rhead{\small 执行准备完成／真实诊断未启动}\cfoot{\small 2026-10-31\quad|\quad\thepage}
\setlength{\parindent}{0pt}\setlength{\parskip}{5pt}
\renewcommand{\arraystretch}{1.18}\setlength{\tabcolsep}{5pt}
\setlist{nosep,leftmargin=1.7em,topsep=4pt}
\ctexset{section={format=\large\bfseries,beforeskip=9pt,afterskip=4pt},subsection={format=\normalsize\bfseries,beforeskip=6pt,afterskip=3pt}}
\emergencystretch=3em
\newcommand{\key}[1]{\textcolor{ink}{\textbf{#1}}}
\newcommand{\code}[1]{{\small\texttt{\detokenize{#1}}}}
\begin{document}
{\LARGE\bfseries 极端降水诊断}\\[4pt]
{\Large 一次执行准备、限额与准入门槛}\\[6pt]
2026年10月31日\quad 版本V1\quad 研究者主导，AI辅助

\section*{1. 此次推进的目的与实际状态}
\key{用一次可追溯的联合诊断，区分点估计转换、条件分位数尾部与观测证据的缺口，帮助研究者决定下一步投入。}
既有A/B/C三种点估计在九个模型的严格$R>20$ mm/h诊断中，真阳性TP=0、假阳性FP=0、漏报FN=717；这不等于全部概率预测失败。保存的最高条件分位数q32还不能说明其高值是否与真实强雨同地同时。

本轮完成的是执行准备：新适配器、限额控制、固定代码副本、脱敏提案与合成检查。没有启动训练、读取气象／标签／静态包载荷、加载权重或运行模型，2025最终测试集继续封存。未生成启动批准文件，也未创建正式运行输出目录。

\begin{center}\small
\begin{tabularx}{\linewidth}{@{}Xl@{}}
\toprule
检查对象 & 当前结果\\
\midrule
数学边界、联合导出、一次性／读取／限额门槛合成检查 & 38项通过\\
九模型历史检查点身份、字节数与开发人口元数据 & 已绑定\\
旧评价器的19份原指标一致性收据 & 继承，未重复运行\\
新入口与独立后台入口无批准时的拒绝测试 & 均通过，未创建启动声明\\
真实联合适配器的端到端数值一致性 & 尚未验证\\
科学验收、强雨风险能力及物理因果解释 & 等待研究者判断\\
\bottomrule
\end{tabularx}
\end{center}

\textbf{固定执行代码：}\\
\code{721bfa979a008e0fc2bbdf031a1888a502ceaca9}\\
\textbf{候选提案规范化SHA-256（安全散列算法256位）身份：}\\
\code{b8ec5e7061c9625715d0a4954dd00cff}\\
\code{707de999cb3a7712166373861ee00de5}\\
以上两行连接为完整提案身份。文档提交可以在该代码之后发生，执行代码固定，不随文档分支头漂移。

\section*{2. 一次性任务范围}
只使用既有B4、B7、B8的2029／2030／2031种子、E0损失、固定第九轮末检查点E9，不选最佳轮或最佳种子，不增加模型。

先做一次2023训练标签／掩码共享统计，再按以下顺序串行处理九份模型：\\
\code{B7/2029, B8/2029, B4/2029, B7/2030, B8/2030, B4/2030,}\\
\code{B7/2031, B8/2031, B4/2031}。

每份模型只做一次2024开发人口完整前向。开发963场景、3,303,090个有效像元时次、有雨407,820个；保持原监督，不因地下压力层缺失而删除监督像元。主比较B8减B7；次要B7减B4、B8减B4，先逐种子配对。

\newpage
\section*{3. 研究者需要冻结的方法}
\textbf{保留原数值路径。}
固定旧评价器执行原模型前向与原指标，旁路只复制已经计算的标签、发生概率、32个条件对数分位数及同一已解码气象帧中的上下文。原概率使用FP32（32位浮点）精确承载已有BF16（脑浮点16位）或FP32取值，不重新计算逻辑函数；分位数保留FP64（64位浮点）。

\textbf{保留联合对应。}
私有输出包含场景、分析时刻、监督像元索引、真值与概率输出，保存分块字节身份。B4地形／水汽分层关联到本包首个B7记录，逐批核对时空键和真值完全一致。不同模型不得各自在不同支持子集上排名。

\textbf{阈值与候选点估计。}
主阈值严格$R>20$，次要5、10、50 mm/h，沿用FP32严格比较；另报恰等于20、50的数量。A/B/C三种方案全部报告，零检出、未定义虚警比、退化与改善均保留，不根据开发结果挑选最终方案。0.5概率门槛仅作诊断参考。

\subsection*{有限分位数支持什么}
发生概率$p=P(R>0.1\mid X)$，条件分布针对$R>0.1$。令
$z=\log(1+r)$，分位水平$\tau_i=(i-0.5)/32$，条件对数分位数$q_i$。
\[
 L_F(z)=\max\bigl(\{0\}\cup\{\tau_i:q_i\le z\}\bigr),\quad
 U_F(z)=\min\bigl(\{1\}\cup\{\tau_i:q_i>z\}\bigr).
\]
\[
 P(R>r\mid X)\in[p(1-U_F),\,p(1-L_F)].
\]
不将未知尾部补零，也不拟合新尾部分布。主概率产出为Brier Score（布里尔分数）的范围、AUROC（受试者工作特征曲线下面积）的含并列排序范围，以及可识别宽度。AP（平均精确率）保留“无法点识别”。差值范围不能确定正负时，明确报告“收益方向尚未识别”；范围不是置信区间。

\subsection*{条件雨强、分层与参考统计}
报告全部32分位覆盖、Conditional Pinball Loss（条件分位数损失）、q32对数欠覆盖幅度、A/B/C误差与阈值检出计数。q32名义水平为$31.5/32=0.984375$；以真值雨强分层后的覆盖仅描述，不能逐层强求名义覆盖。

保留七个真值雨强层$(0.1,1],(1,5],(5,10],(10,20],(20,30],(30,50],(50,+\infty)$、四个UTC（协调世界时）时次，以及低于1000、1000--2000、不低于2000米高程。DOTE（动态地形与热力／水汽强迫表征）按850／700百帕方向符号、压力可用性与水汽交叉分层；水汽参考固定31.85768219039871 kg/m$^2$，不由开发指标选择。

共享训练标签参考使用同人口2023的超阈频率，以及$R>0$的P95／P99（第95／99百分位），$R>0.1$仅作敏感性口径，样本分位采用线性定义。无新观测和独立事件分组；717个强雨像元时次仍不是717场暴雨。

\newpage
\section*{4. 拟议资源限额：当前尚未获批}
\begin{center}\small
\begin{tabularx}{\linewidth}{@{}Xr@{}}
\toprule
项目 & 一次包限额／固定参考\\
\midrule
九次开发气象／掩码读取 & 239,753,278,008字节／71,550内容实例\\
六次静态包读取 & 1,876,968字节／6实例\\
一次共享训练标签／掩码 & 180,049,094字节／246实例\\
上述合计 & 239,935,204,070字节／71,802实例\\
九检查点，各散列一次、反序列化一次 & 964,353,774档案长度参考字节／18打开\\
载荷及检查点档案长度参考合计 & 240,899,557,844字节\\
受控工作进程冻结元数据档案预留 & 至多1,073,741,824字节／2,048打开\\
\textbf{上述源预留总上限} & \textbf{241,973,299,668字节}\\
\midrule
总时限，取墙钟与单调时钟更早期限 & 21,600秒（6小时）\\
单模型上限／共享标签上限 & 2,400秒／600秒\\
每工作进程元数据上限 & 64 MiB／128打开\\
新输出总容量上限 & 16 GiB\\
启动及每任务前最低空闲磁盘／主机内存 & 32 GiB／8 GiB\\
最低空闲GPU（图形处理器）显存 & 5,120 MiB\\
优化器更新／重试／重放 & 0／0／0\\
\bottomrule
\end{tabularx}
\end{center}

\textbf{口径限制：}字节为逻辑源交付及打开前档案长度预留，不能宣称为操作系统物理磁盘读取量。检查点内部原生压缩包访问不冒充逐字节实测。固定代码、库与批准身份检查单独绑定；表中元数据上限针对受控工作进程内的冻结元数据打开。不存在6小时之外另加标签时间的额度。

\textbf{当前资源检查：}磁盘约47.5 GiB、可用内存约14.9 GiB、空闲显存约6.3 GiB，快照满足上述准入线。资源会变化；正式启动和每任务前重新检查。失败不自动放宽限额，不终止其他研究任务。

\section*{5. 验收与停止规则}
\begin{enumerate}
\item 新批准须同时绑定研究者决策编号、明确指令、执行代码、提案身份、全部方法、任务顺序、预算与全新输出目录。历史批准不复用；本文件不构成批准。
\item 权重校验与严格加载通过后，原评价器的保存指标按绝对／相对$10^{-12}$容差及精确整数计数核对。每模型必须963场景、3,303,090监督、407,820有雨及7,950源内容实例、26,639,253,112交付字节。
\item 联合记录须时空键唯一、顺序一致、监督一致；跨模型分层必须像元成员散列一致。原指标门槛或任何身份、时限、容量、读取计数失败，立即停止相关后续任务，保留不确定读取和所有部分输出。
\item 仅全部工程门槛通过才封口。真实诊断完成后状态是“等待研究者科学审查”，不自动批准模型、物理机制、独立事件能力或论文结论。
\end{enumerate}

\key{现在仅需研究者集中决定：是否按上述固定方法、代码与限额，批准该一次性诊断包。}
批准不包含后续训练、独立观测采集、B9、2025访问或失败后的再次启动。

{\footnotesize
可追溯材料：同目录的proposal.json、engineering\_verification.json、preparation\_receipt.json及可修改main.tex。\\
代码：\href{https://github.com/flyeagle-cex/YunTAPR-Net/commit/721bfa979a008e0fc2bbdf031a1888a502ceaca9}{721bfa9固定执行版本}。
旧证据基线7c19b9d，上一阶段准备923f250；本文件没有替代其历史收据。}
\end{document}
